% =====================================================================
%  Chapter 12 — The State of Play, Part III
% =====================================================================

\chapter{The Review That Has Not Happened}
\label{ch:noreview}

The most consequential fact about the claim under examination is an absence,
and absences are hard to write about without either overstating or
understating them. This chapter attempts to state it precisely.

\section{What exists}\begin{itemize}
\item \textbf{The remark.} Four pages, at \path{alpo.ge/remark.pdf}, dated by
its mirror to 20 November 2026. No theorem statement, no proof, no
definitions.
\item \textbf{The Lean development.} 239{,}275{,}554 bytes, 175{,}187 files,
roughly 17.4 million lines, pinned to Lean \texttt{v4.33.1} and Mathlib
\code{0df444a360eaa60ab8c11dca51a86af692955474}.
\item \textbf{A third-party mirror} with a README, a static scan, and a set of
recommended verification steps.
\item \textbf{Two independent compilations}, reported the same day.
\item \textbf{Press coverage} in Quanta Magazine, New Scientist, Scientific
American, \emph{The Economist}, and a wire report from Agence France-Presse.
\end{itemize}

That is a complete and, by the standards of announcements, admirably
well-organised public record.

\section{What does not exist}\begin{itemize}
\item \textbf{No preprint.} No arXiv paper, no journal submission, no
technical report describing the argument. A search of arXiv returns exactly
one item connected to this story --- arXiv:2608.13637 --- and it is about the
two-thirds theorem, not about RH.
\item \textbf{No refereed report.}\idx{refereed report, absence of} No expert has published an assessment.
\item \textbf{No statement review.} The human check described in Chapter~6
--- reading the formal statements against the classical theorems they claim to
encode --- has not been performed by anybody outside the author.
\item \textbf{No written proof.} The remark's author says the prose version
``is routine'' to produce but has not been written, and the author's own
summary is what he calls ``slop''.
\end{itemize}

The wire report put the position plainly:

\begin{quote}
\small
``The claim has not been independently verified, and the proof has not yet been
reviewed by outside experts.''
\end{quote}

\section{Why this absence is the central fact, and not a technicality}A formalised proof of the Riemann hypothesis would be the most consequential
mathematical result of the century, and it would be the most consequential
\emph{proof}. Unlike a paper, it could in principle be checked by anybody with
a computer. That is the promise of formalisation, and this is the first
candidate of that kind of scale.

The absence of a review is not a bureaucratic delay. It is the substantive
step. Consider what a referee must do with this artifact:

\begin{enumerate}
\item Build it. Done, twice.
\item Check the axiom list. Done, mechanically.
\item Read the eight-link chain. Done in this book, for this book.
\item \textbf{Read the statements of the load-bearing inputs --- the theta
lift, the automorphy bundle, the admissibility condition, the vanishing
theorem --- against the published theorems they encode.}
\item \textbf{Read the applications: check that the classical theorems are
applied in settings where their hypotheses actually hold.}
\item \textbf{Compare the formalisation with the remark's own account, and
resolve the 15-to-16 November gap.}
\end{enumerate}

Steps 1--3 are the part that has been done. Steps 4--6 are the part that has
not, and they are the part that could go either way.

\begin{seamnote}{On the scale of the work}
Step 4 alone is not a single afternoon's reading. The inputs are drawn from
the theory of automorphic forms, the theory of Galois representations, the
theory of Iwasawa algebras, and the theory of theta lifts --- four substantial
fields, none of them easy. The relevant classical statements must be located
in the literature, compared field by field with the Lean definitions, and the
comparison written up.

The realistic estimate for a competent specialist is weeks, not days --- and
that is for one person, reading one input. There are five inputs named in
\tabref{tab:inputs} and a long tail behind them.

This is why the community's normal mechanism for this kind of work --- a
specialist reads the development and reports --- has not produced a report
yet. It is not that nobody has looked. It is that the job is large and the
incentive structure of the field does not make it anyone's clear priority
until the claim is on the record in a form that can be cited.
\end{seamnote}

\section{What the absence does not mean}It is important not to over-read the absence, because the available evidence
cuts in both directions and the temptation to pick a side is strong.

\begin{quote}
\textbf{It is not evidence of a defect.} No referee has reported a defect. No
mathematician has published a refutation. The press coverage contains
scepticism, and one prominent response --- Umberto Borretti's ``Re: Riemann''
--- but a sceptical essay is not a refutation, and the absence of a refutation
is not a refutation of the absence.

\textbf{It is not evidence of a good proof either.} ``Nobody has found a
problem'' is not a positive statement about a development of 17.4 million
lines. Nobody has read it.

\textbf{It is evidence about the normal process, not about the mathematics.}
What it tells us is that the claim has not yet passed through the mechanism
that would have made it a settled result. That mechanism is slow, and it is
the same mechanism that every other proof claim has gone through.
\end{quote}

\section{The precedent that matters}Claims to have proved the Riemann hypothesis have a long history, and the
usual outcome is a long, careful, unglamorous process of elimination. This
book does not need the parade of earlier claimants, most of which are
low-quality; their existence is well known and it establishes only that the
conjecture attracts attempts.

The more useful precedents are two, and they point in opposite directions.

\begin{description}[style=nextline,leftmargin=0pt,labelwidth=0pt,labelsep=0pt,font=\sffamily\bfseries\small,
  itemsep=5pt]
\item[The successful formalisations.] Several substantial results in number
theory and combinatorics have been formalised in Lean, and a few have been
formally verified at a level that the community accepts: the
\textbf{Prime Number Theorem for arithmetic progressions}\idx{prime number theorem, formalised}, the
\textbf{Four Colour Theorem}\idx{four-colour theorem}, the \textbf{Boolean Prime Ideal Theorem},
\textbf{Apéry's theorem}, and others. The pattern is consistent: the
formalisation is accompanied by a paper describing the mathematics, a
statement review, and an independent check of the kernel. Those results are
uncontroversial because they went through the process.

\item[The withdrawal.] OpenAI's own project history records a withdrawal in
October 2026 with a stated sign error. The development compiled. The axiom
list was clean. The error was mathematical. This is the specific scenario
that a review catches and a compiler does not.
\end{description}

Taken together, the precedents say two things. They say that the machinery
works, and that a formalised proof of RH is not an absurd thing to attempt.
And they say that the specific gap between ``compiles'' and ``accepted'' is
real, routinely crossed, and routinely not.

\section{Why the author's own words matter here}One of the most striking features of the remark, reading it as an
evidentiary document rather than a news story, is how consistently it
disclaims the kind of authority that would make a review unnecessary.

\begin{quote}
\small
``in my view of my path I have just taken a step'' \qquad
``which might make zero sense, and I apologize, though the Lean should suffice
to clarify everything. I'm just too inexpert in the relevant stuff at the
moment'' \qquad
``I personally have little mastery of automorphic forms on higher rank groups
(youthful foolish stubbornness)'' \qquad
``I also emphasize that I may have baked lots of strange suboptimal choices
into the argument''
\end{quote}

These are not the words of a person claiming that review is a formality. They
are the words of a person who knows that review is the substance, and who has
put the work in a form where a reviewer can do it --- because the author cannot
do it himself, and says so.

There is a reading of the remark in which this is the strongest possible
stance: the author has produced a machine-checkable artifact precisely so that
the part he cannot do is done by others. The remark's ``I decided to only
present Lean'' is then not a retreat from scrutiny but a \emph{design} for it:
the best available evidence, in the most checkable form, with the author's
honest assessment of its state attached.

That reading is available. So is the other one, in which the same document
amounts to an assertion with no supporting argument. The evidence currently
available does not decide between them, and the difference between them is
exactly what a review would establish.

\begin{keytakeaways}

\item There is no preprint, no refereed report, no statement review and no
written proof. There is a remark, a mirror, two compilations, and press
coverage.
\item The steps that have been completed are the mechanical ones: build, axiom
check, and the structural trace of Chapter~7. The steps that remain are the
mathematical ones, and they are the ones that decide the question.
\item The absence of a review is evidence about the process, not about the
mathematics. It is not evidence of a defect and it is not evidence of
correctness.
\item The remark itself consistently disclaims the authority that would make
review unnecessary. That disclaimer is the most credible thing in the
document, and it is also the reason a review is needed.
\end{keytakeaways}
